\section{Materials and Methods}

\subsection{Patient tumour and cell line RNA-seq data collection}

RNA-seq data were assembled for DSMZ cancer cell lines and public patient tumour cohorts. Raw RNA-seq FASTQ files from 167 cancer cell-line RNA-seq profiles were obtained from the Leibniz Institute DSMZ (Deutsche Sammlung von Mikroorganismen und Zellkulturen GmbH; German Collection of Microorganisms and Cell Cultures). The panel included breast cancer, neuroblastoma, retinoblastoma, and haematological malignancy profiles. Patient tumour RNA-seq count data for TCGA cohorts were obtained through the Genomic Data Commons (GDC) portal \citep{GeistlingerRamos2025}. Raw RNA-seq FASTQ files for additional patient tumour cohorts were obtained from the Gene Expression Omnibus (GEO) or Sequence Read Archive where applicable. Controlled access sequencing data are not redistributed.

Primary tumour--cell-line modelling and bidirectional similarity-ranking analyses were restricted to breast cancer, neuroblastoma, and retinoblastoma. Haematological malignancy cell-line profiles were retained only as an auxiliary comparator lineage for cell-line-only pan-cancer analyses where explicitly stated. Haematological patient cohorts were catalogued but were not used in the primary tumour--cell-line modelling or bidirectional ranking analyses.

\begin{table}[H]
\centering
\caption{Overview of RNA-seq datasets processed or catalogued for the study.}
\label{tab:dataset_overview}
\begin{tabular}{lcc}
\toprule
Cancer type & Cell line RNA-seq profiles & Patient tumour samples \\
\midrule
Breast cancer & 29 & 1231 \\
Neuroblastoma & 18 & 282 \\
Retinoblastoma & 11 & 66 \\
Haematological malignancies & 109 & 223 \\
\bottomrule
\end{tabular}  
\end{table}


\begin{table}[H]
\centering
\caption{Public RNA-seq datasets used for patient tumour samples obtained from TCGA and GEO.}
\label{tab:patient_datasets}
\begin{tabular}{lll}
\toprule
Cancer type & Dataset accession & Samples \\
\midrule

Breast invasive carcinoma & TCGA-BRCA & 1231 \\

Neuroblastoma & GSE100148 & 4 \\
Neuroblastoma & GSE189367 & 64 \\
Neuroblastoma & GSE49711 & 5 \\
Neuroblastoma & SRP409177 & 47 \\
Neuroblastoma & TARGET-NBL & 162 \\

Retinoblastoma & GSE111168 & 6 \\
Retinoblastoma & GSE196420 & 18 \\
Retinoblastoma & GSE268136 & 42 \\

Burkitt lymphoma & GSE199413 & 8 \\
Chronic lymphocytic leukaemia & GSE237907 & 10 \\
Diffuse large B-cell lymphoma & TCGA-DLBC & 48 \\
Acute myeloid leukaemia & TCGA-LAML & 151 \\
Myeloproliferative neoplasms & GSE228758 & 6 \\

\bottomrule
\end{tabular}
\end{table}


\subsection{Study design and units of analysis}

The study used a patient-referenced transcriptomic design in which patient tumour cohorts and cancer cell-line RNA-seq profiles were analysed in shared expression spaces after source harmonisation. Different analyses used different statistical units, and the unit was matched to the corresponding question.

An RNA-seq library/profile denotes one processed sequencing profile from one cell-line library or one patient tumour sample. A biological cell-line group denotes a canonical cell-line identity after grouping replicate libraries that represent the same cell line. A patient tumour sample denotes one tumour RNA-seq profile from TCGA, GEO, TARGET, or SRA-derived cohorts. A feature--distance representation denotes one selected gene set paired with one distance metric or transformed expression space. A cell-line--tumour similarity pair denotes one comparison between a cell-line profile or biological cell-line group and one patient tumour sample. A graph node denotes the entity represented in the graph under analysis, and an edge denotes a retained similarity relationship under the corresponding graph rule. A marker-gene contrast denotes one focal cell-line profile, biological cell-line group, or graph-derived anchor compared with its lineage-matched REST group. An enrichment gene set denotes the marker list passed to g:Profiler together with its eligible background universe.

Patient-to-cell-line ranking retained replicate libraries as profile-level cell-line models. This direction therefore used 58 solid-tumour cell-line profiles: 29 breast cancer profiles, 18 neuroblastoma profiles, and 11 retinoblastoma profiles. RBL\_15 and RBL\_20 each contributed two replicate libraries in this profile-level analysis.

Cell-line-centred tumour-class similarity and Precision@k analyses collapsed replicate libraries to biological cell-line groups before ranking patient tumour samples. The primary group-level evaluation therefore used 56 biological cell-line groups: 29 breast cancer groups, 18 neuroblastoma groups, and 9 retinoblastoma groups. The biological-group-level retinoblastoma count is lower than the profile-level retinoblastoma count because the replicate libraries for RBL\_15 and RBL\_20 were each collapsed to one biological group.

Graph analyses used the node set specified for the corresponding graph. Patient-referenced per-lineage graphs and resolved neighbour graphs used cell-line profiles or profile-normalised cell-line identifiers according to the lineage-specific input tables, whereas the pan-cancer cell-line network retained the 167 analysed RNA-seq profiles as profile-level nodes. Consequently, retinoblastoma counts can differ across profile-level graphs, biological-group ranking analyses, and pan-cancer profile-level networks.

\subsubsection{Threshold rationale}

Numerical thresholds were used as pipeline parameters or fixed analysis criteria rather than as a single formal hypothesis-testing framework. The tumour-purity cutoff, feature-retention sizes, cluster-number grid, consensus-support cutoff, graph sparsification rules, nearest-neighbour depth, top-\(k\) reporting cutoffs, marker-gene filters, and enrichment significance threshold were encoded in the pipeline configuration or fixed analysis scripts. They were selected for interpretability, minimum sample support, computational tractability, or established practice for the corresponding method. These thresholds were exploratory and descriptive rather than pre-registered, and they were not chosen to optimise the final lineage-classification or tumour-class similarity outcomes.



\subsection{RNA-seq data cleaning and processing}

\subsubsection{Preprocessing steps and quality control}

Raw paired-end RNA-seq FASTQ files from the analysed cell-line panel and non-TCGA patient tumour cohorts were processed using a workflow analogous to the GDC mRNA-seq pipeline \citep{GDCpipeline}. Sequencing quality assessment was performed using FastQC (v0.11.9). Reads were aligned to the human reference genome (GRCh38/hg38) using STAR (v2.7.10b) in two-pass mode \citep{Dobin2013}. The resulting alignments were coordinate-sorted and indexed with SAMtools (v1.17) \citep{Danecek2021}. Library strandedness was inferred from aligned BAM files using RSeQC (v5.0.1) \citep{Wang2012}. Transcriptome-wide read coverage was computed with a custom Python script implemented using \texttt{pysam} (v0.21.0). Gene annotations were obtained from GENCODE~v44 (GRCh38; Ensembl release~110). Alignment-level quality metrics were collected for each BAM file using Picard tools (v3.0+) \citep{Picard2019}. Project-level quality-control reports were generated with MultiQC (v1.14), which aggregated FastQC, STAR, and Picard outputs \citep{Ewels2016}. TCGA patient tumour cohorts were not realigned; GDC gene-level count matrices were harmonised for subsequent analyses by Ensembl gene identifier.

\subsubsection{Tumour purity estimation}

Tumour purity was estimated for patient tumour samples using the Estimation of STromal and Immune cells in MAlignant Tumours using Expression data (ESTIMATE) algorithm \citep{Yoshihara2013}. ESTIMATE infers stromal and immune infiltration from expression signatures. To minimise confounding from non-tumour cellular admixture, patient tumour samples were required to have an estimated tumour-purity score of $\geq 0.7$ for inclusion in tumour--cell-line analyses. The threshold of 0.7 was selected with reference to the ESTIMATE validation framework and to prior TCGA transcriptomic analyses using similar cutoffs to reduce stromal and immune admixture \citep{Jin2023}. Pearson and Spearman correlations were calculated between tumour-purity scores and immune or stromal scores for each data source as quality-control statistics.

After purity filtering, analysis-specific tumour sets were additionally required to be present in the harmonised expression matrix, to have non-missing source and lineage metadata, and to contain the genes required for the corresponding feature space. Consequently, the purity-retained sample count and the sample count used in a specific feature-selection, graph, or ranking analysis may differ; analysis-specific sample sizes are reported with the corresponding figures and tables.

\subsubsection{Source harmonisation and normalisation}

Gene-level count matrices from cancer cell lines and patient tumour samples
were harmonised by matching shared gene identifiers, using Ensembl gene IDs,
and retaining the intersection of genes present in both datasets. For each
disease lineage, patient tumour samples and cell lines were combined into a
shared count matrix.

RNA-seq count data are discrete and typically exhibit overdispersion, whereby
the variance increases with the mean. Unwanted technical variation across
datasets can introduce systematic differences between samples processed under
different experimental conditions \citep{Leek2010BatchEffects,Leek2012}. Source
harmonisation was therefore performed using ComBat-seq, which models RNA-seq
count data using a negative binomial framework \citep{Zhang2020ComBatSeq}.

ComBat-seq was applied separately within each disease lineage, with sample
source specified as the batch variable. The source variable encoded the
dataset or broad origin of each sample, including cell-line and patient tumour
sources as represented in the lineage-specific input metadata. Because
harmonisation was performed within disease lineage, disease lineage was not
included as a covariate in the ComBat-seq model. No biological cell-line group,
tumour class, graph-derived label, similarity-ranking outcome, or
differential-expression result was included as a covariate. This approach was
chosen because it preserves count scale before variance stabilisation and
provides a count-aware source harmonisation step for mixed RNA-seq data sources.

Source can be partially confounded with biological sample type, because
patient tumours and cell lines originate from different collection and
processing contexts. The ComBat-seq step was therefore treated as source
harmonisation rather than complete batch removal. Unsupervised structure was
inspected after correction with UMAP and clustering outputs, but source
harmonisation does not guarantee removal of all unwanted variation or
preservation of all biological signal.

Source-harmonised count matrices were subsequently transformed using the
variance stabilising transformation implemented in \texttt{DESeq2}
\citep{Love2014}. VST reduces the dependence of variance on the mean, producing
approximately homoscedastic data suitable for clustering and distance-based
analyses.

UMAP was used exclusively as a qualitative visualisation of transcriptomic
structure before and after source harmonisation \citep{McInnes2020}. No quantitative
inference was made from UMAP coordinates.


\subsection{Feature selection methods}
\subsubsection{Unsupervised feature selection}

Unsupervised feature selection was performed to define expression features without using sample class labels \citep{Solorio-Fernandez2020}. Each feature-selection criterion, combined with a distance metric, defined one feature--distance representation. The procedure consisted of three stages (Supplementary
Algorithm~\ref{alg:supp_unsupervised_feature_selection}): (i) pre-filtering of genes with near-constant expression values, (ii) first-stage statistical ranking to define a candidate gene pool, and (iii) network-based refinement.

\paragraph{Pre-filtering.}
Genes with near-constant expression across the variance-stabilised expression matrix were removed by excluding genes with standard deviation $\leq 10^{-8}$. 

\paragraph{First stage statistical ranking.}
To quantify complementary aspects of transcriptional variability, six criteria
for unsupervised gene ranking in the first stage were applied independently.
Genes were ranked by variance, highly variable gene (HVG) residual variance,
median absolute deviation (MAD), mean absolute deviation (MeanAbsDev), Shannon
entropy, and scores derived from principal component analysis (PCA) loadings.
HVG scores were computed as residuals of log variance after adjustment for the
mean--variance trend. For each
gene, the score was defined as the observed \(\log_{10}\) variance minus the
expected \(\log_{10}\) variance at the same mean expression level, estimated
using a LOESS fit. Shannon entropy was calculated after discretising each
gene's expression values into 20 equal width bins. PCA based feature selection
was performed on the mean centred expression matrix, and genes were ranked by
the summed absolute loadings across the first five principal components, or
fewer when fewer than five principal components were available. Ties for all
first-stage rankings were resolved by Ensembl gene identifier.

For each criterion, the top 3000 genes were retained. The union of these six
sets defined the candidate gene pool used for subsequent network-based
refinement.

\paragraph{Network-based refinement.}

To prioritise genes embedded in coherent co-expression structure, network-based metrics were computed on the candidate gene pool.

First, a Spearman connectivity score was calculated for each gene as the mean absolute Spearman correlation to all other genes within the candidate set.

Second, the MX score was defined as a combined connectivity and variance score that prioritised genes with high mean absolute Spearman connectivity and high variance:
\[
\mathrm{MX}(g)=K_S(g)\times \frac{\mathrm{Var}(g)}{\max_h \mathrm{Var}(h)},
\]
where \(K_S(g)\) denotes the mean absolute Spearman connectivity of gene \(g\) within the candidate set.

Third, weighted gene co-expression network analysis (WGCNA) was used to compute total network connectivity \citep{Langfelder2008WGCNA}. Candidate-gene expression values were transposed to sample-by-gene format and standardised per gene. Soft-threshold powers \(1{:}10,12,14,\ldots,30\) were evaluated with \texttt{WGCNA::pickSoftThreshold}; the selected power was the candidate with the highest scale-free topology fit among powers with \(R^2 \geq 0.8\) and slope between \(-2.5\) and \(-1.5\), with power 5 used when no candidate satisfied both criteria. After soft-threshold selection, an unsigned adjacency matrix was constructed. Total network connectivity (\textit{kTotal}) was then computed as the sum of adjacency weights linking each gene to all other genes in the network.

For each network-based metric (Spearman connectivity, MX score, and WGCNA \textit{kTotal}), genes were ranked in decreasing score order. Genes with identical scores were ordered lexicographically by Ensembl gene identifier to make the ranking deterministic, and the top 500 genes were retained.

Together, the six first-stage statistical criteria and three network-based metrics defined nine unsupervised feature selection methods. For breast cancer analyses, the PAM50 gene signature was additionally assessed as a predefined breast cancer feature set using both Euclidean and correlation distance representations; PAM50 was not used as an unsupervised feature selection method for other cohorts.


\begin{figure}[H]
    \centering
    \includegraphics[width=0.8\textwidth]{Thesis_writeup/Figures/unsupervised_feature_selection_pipeline_v3.pdf}
    \caption{Overview of the unsupervised feature selection pipeline. In Phase 1, six complementary statistical criteria (Variance, HVG, MAD, MeanAbsDev, Entropy, and PCA loadings) independently rank genes and retain the top 3,000 each. The union of these rankings defines a candidate pool that serves as input to Phase 2, in which three network-based methods (Spearman connectivity, MX score, and WGCNA kTotal) identify the top 500 genes under each network metric. All nine resulting feature sets are propagated independently for clustering so that feature-set-specific representations remain separate.}
    \label{fig:feature_selection_pipeline}
\end{figure}

\subsubsection{Clustering analysis}

Following unsupervised feature selection, clustering was performed on the
VST-transformed expression matrix containing both patient tumour samples and
cell lines, restricted to the selected feature sets (Supplementary
Algorithm~\ref{alg:supp_clustering_consensus}). Analyses were conducted in two
feature spaces: (i) the expression matrix restricted to the selected features
and (ii) its top 20 principal components.

The number of clusters was evaluated over the configured predefined range,
\[
k \in \{2, \ldots, 8\},
\]
excluding values of \(k\) that exceeded the number of available samples or
informative features. Cluster quality was assessed using the mean silhouette width
\citep{Rousseeuw1987Silhouette}.
In the case of ties, the smaller
value of \(k\) was chosen; remaining ties were resolved deterministically by
representation identifier.

\paragraph{Hierarchical clustering.}
Agglomerative hierarchical clustering was performed using Ward's minimum variance method (Ward D2), with pairwise sample dissimilarities quantified using Euclidean distance and Pearson correlation-based distance \((1-\rho)\) \citep{BertholdHoeppner2006}.

\paragraph{K-means clustering.}
K-means clustering was performed in Euclidean space with 20 random initialisations. The solution achieving the lowest within-cluster sum of squares was selected \citep{MacQueen1967}.

\paragraph{Consensus clustering.}
Cluster robustness was assessed using consensus clustering \citep{Monti2003ConsensusClustering}, implemented with \texttt{ConsensusClusterPlus} \citep{Wilkerson2010ConsensusClusterPlus}, with 1{,}000 resampling iterations, retaining 80\% of samples and 100\% of features per iteration. Hierarchical clustering used the relevant distance matrix, whereas k-means was run only for Euclidean-compatible representations. To accommodate the internal Euclidean distance implementation of consensus clustering, correlation-based analyses were performed after mean-centring and unit-\(\ell_2\)-normalising samples, such that Euclidean distance in the transformed space corresponded to Pearson correlation-based distance in the original space. The optimal number of clusters was selected as the value of \(k\) that maximised the change in the area under the cumulative distribution function of the consensus matrix (\(\Delta\)AUC).

\subsubsection{Tumour neighbourhood analysis}

Tumour neighbourhood analysis was performed to identify patient tumour
samples most similar to each cancer cell line within the corresponding
joint expression analysis of cell lines and tumour samples (Supplementary
Algorithm~\ref{alg:supp_tumour_neighbourhood_inference}). For each
feature--distance representation \(q\), the input matrix
\(\mathbf{X}^{(q)} \in \mathbb{R}^{N_q \times G_q}\) contained
variance-stabilised expression values for \(N_q\) samples and
\(G_q\) selected genes, where \(N_q\) denotes the number of cell-line and
patient tumour samples included for representation \(q\), and \(G_q\) denotes
the number of genes assigned to that representation. Rows represented samples and
columns represented selected genes. For cell-line-centred neighbourhood
inference, redundant technical sequencing identifiers were collapsed after
transformation to the canonical profile identifier used by the corresponding
analysis by averaging transformed expression values. Cluster labels for
collapsed identifiers were assigned by majority vote with deterministic
tie-breaking by cluster label. This profile-normalisation step was distinct
from the biological-group collapse used for cell-line-centred tumour-class
similarity and Precision@k.

Cluster assignments from the corresponding upstream clustering output were
mapped to the sample identifiers used in the neighbourhood analysis.
Neighbourhood analysis was restricted to samples present in both the expression
matrix and the cluster assignment vector.

For each feature--distance representation \(q\), tumour neighbourhoods were computed
separately for each implemented clustering method \(m\). The method set
included k-means, hierarchical clustering, and consensus clustering outputs.
Hierarchical clustering was implemented for both Euclidean and correlation
distance measures, whereas k-means was implemented for the Euclidean distance
measure. Let
\(\mathcal{S}_{qm}\) denote the set of cell lines,
\(\mathcal{P}_{qm}\) the set of patient tumour samples, and
\(c_i^{(qm)}\) the cluster assignment of sample \(i\) under method
\(m\). For a focal cell line \(s \in \mathcal{S}_{qm}\), the initial
candidate tumour set was defined as the patient tumour samples
assigned to the same cluster as the focal cell line:

\begin{equation}
\widetilde{T}_{qm}(s) =
\{ t \in \mathcal{P}_{qm} : c_t^{(qm)} = c_s^{(qm)} \}.
\end{equation}

If fewer than 10 patient tumour samples were assigned to the same
cluster as the focal cell line, the candidate set was expanded to all
patient tumour samples included for that method:

\begin{equation}
T_{qm}(s) =
\begin{cases}
\widetilde{T}_{qm}(s), & |\widetilde{T}_{qm}(s)| \geq 10, \\
\mathcal{P}_{qm}, & |\widetilde{T}_{qm}(s)| < 10 .
\end{cases}
\end{equation}

Patient tumour samples in \(T_{qm}(s)\) were ranked by increasing
distance to the focal cell line, with the metric determined by the
representation label. For Euclidean representations, patient tumour samples
were ranked by Euclidean distance in the selected expression space. For
correlation representations, samples were ranked by Pearson correlation
distance, defined as one minus the Pearson correlation coefficient
between the selected gene expression profiles of the cell line and the
patient tumour sample. Exact distance ties were resolved deterministically by patient tumour sample identifier. The ordered candidate list for cell line \(s\)
was denoted
\((t_{(1)}, t_{(2)}, \ldots, t_{(|T_{qm}(s)|)})\).

Neighbourhood size was determined adaptively from the ranked candidate
set. For each cell line, the selected neighbourhood contained the top
10\% of candidate patient tumour samples, constrained to a minimum
of 30 and a maximum of 200 samples, without exceeding the number of
available candidates:

\begin{equation}
k_{qm}(s) =
\min\!\left(
\max\!\left(30,
\left\lceil 0.10 \cdot |T_{qm}(s)| \right\rceil
\right),
200,
|T_{qm}(s)|
\right).
\end{equation}

The tumour neighbourhood of cell line \(s\) under method \(m\) was
therefore defined as

\begin{equation}
\mathcal{N}_{qm}(s) =
\{ t_{(1)}, t_{(2)}, \ldots, t_{(k_{qm}(s))} \}.
\end{equation}

\subsubsection{Consensus aggregation}

Consensus aggregation was applied to quantify the stability of neighbourhood
assignments from cell lines to patient tumour samples across clustering
methods within each feature--distance representation (Supplementary
Algorithm~\ref{alg:supp_within_direction_consensus}). Within
representation consensus refers to agreement among tumour neighbourhood
assignments computed within the same feature--distance representation. For a
given feature--distance representation \(q\), let \(\mathcal{M}_q\) denote
the set of available tumour neighbourhood result tables for that
representation, and let \(M_q = |\mathcal{M}_q|\). Only discovered
neighbourhood tables contributed to the denominator for that representation.

For each cell line \(s\) and patient tumour sample \(t\), the
number of clustering methods that assigned \(t\) to the selected
neighbourhood of \(s\) was computed as

\begin{equation}
n_{\mathrm{methods}}^{(q)}(s,t) =
\left|
\left\{
m \in \mathcal{M}_q :
t \in \mathcal{N}_{qm}(s)
\right\}
\right|.
\end{equation}

The consensus probability for the pair \((s,t)\) was defined as

\begin{equation}
p_{\mathrm{consensus}}^{(q)}(s,t) =
\frac{
n_{\mathrm{methods}}^{(q)}(s,t)
}{
M_q
}.
\end{equation}

Thus, \(p_{\mathrm{consensus}}^{(q)}(s,t)\) ranged from 0 to 1, where
0 indicated that the pair was not retained by any contributing method
and 1 indicated that the pair was retained by all contributing methods
for that representation. Associations with
\(p_{\mathrm{consensus}}^{(q)} \geq 0.7\) were classified as strong
consensus associations.
Pairs not recovered by any method had implicit \(p_{\mathrm{consensus}}=0\) and were not included in the final consensus table; metrics were computed over the observed neighbourhood union unless otherwise stated.

The threshold \(p_{\mathrm{consensus}} \geq 0.7\) was used as an analysis-specific selection rule to identify strong consensus associations, not as a general statistical significance threshold.

Cell line level metrics were computed over the union of patient tumour
samples assigned to the cell line by at least one contributing
method,

\begin{equation}
U_q(s) =
\bigcup_{m \in \mathcal{M}_q} \mathcal{N}_{qm}(s).
\end{equation}

For each cell line \(s\), the maximum consensus, median consensus, and
fraction of strong consensus associations were computed as

\begin{equation}
\max_{t \in U_q(s)}
p_{\mathrm{consensus}}^{(q)}(s,t),
\end{equation}

\begin{equation}
\mathrm{median}_{t \in U_q(s)}
p_{\mathrm{consensus}}^{(q)}(s,t),
\end{equation}

and

\begin{equation}
\mathrm{frac}_{\geq 0.7}^{(q)}(s) =
\frac{
\left|
\left\{
t \in U_q(s) :
p_{\mathrm{consensus}}^{(q)}(s,t) \geq 0.7
\right\}
\right|
}{
|U_q(s)|
}.
\end{equation}

These metrics were used to compare the stability of tumour
neighbourhood assignments across clustering methods.

\subsubsection{Cross-representation consensus aggregation and stability measures}

Cross-representation consensus measures were computed to quantify the stability
of tumour-neighbourhood assignments across feature--distance representations
(Supplementary Algorithm~\ref{alg:supp_cross_direction_cpi}). For each
cell-line--representation pair, the metrics were the maximum
consensus probability, median consensus probability, fraction of associations
with \(p_{\mathrm{consensus}} \geq 0.7\), and fraction of associations with exact
consensus \(p_{\mathrm{consensus}}=1\). Metrics with near-zero variance were
removed before standardisation.

All four metrics entered CPI in an orientation in which larger values denoted
stronger or more stable neighbourhood support; no metric required sign reversal. The
composite prioritisation index (CPI) was computed from the first principal
component of the standardised metric matrix. Absolute first-component loadings
were normalised to define non-negative metric weights, because the sign of a
principal component is arbitrary whereas loading magnitude records how strongly
each metric contributes to the component. CPI was then calculated as the
weighted sum of the standardised metrics for each cell-line--representation
pair. Larger CPI values indicate stronger relative support across the four
consensus metrics within the analysed set of cell-line--representation pairs.
CPI is a descriptive stability statistic, not a probability, validation metric,
or inferential test. It was not used to construct the pan-cancer feature set,
weight expression profiles, define differential-expression contrasts, or
calculate the final correlation-based similarity rankings.

\subsubsection{Patient-referenced cell-line similarity graphs}

For each feature--distance representation \(r\in\mathcal{R}\), a graph
\(G_r=(S,E_r)\), called the representation-specific graph for \(r\), was constructed
from patient-referenced neighbourhood consensus profiles (Supplementary
Algorithm~\ref{alg:cell_line_similarity_network}).
Here, \(S\) denotes the set of analysed cell lines and
\(E_r\) denotes the edge set retained under representation \(r\).
Each cell line \(s \in S\) was represented as a vector embedded in a
feature space referenced to patient tumour samples
$\mathbb{R}^{P}$, where $P$ denotes the total number of
patient tumour samples. Patient tumour samples absent from the observed
neighbourhood union for a given cell line were assigned consensus probability
zero. The embedding was defined as:

\begin{equation}
\mathbf{v}_s =
\left(
p_{\mathrm{consensus}}(s,t_1), \ldots,
p_{\mathrm{consensus}}(s,t_P)
\right),
\end{equation}

where $t_1, \ldots, t_P$ index the patient tumour samples.

Pairwise similarity between cell lines $s_i$ and $s_j$
was computed using Pearson correlation across the consensus vectors
referenced to patient tumour samples.
This procedure resulted in a symmetric similarity matrix

\begin{equation}
\mathbf{S} \in [-1,1]^{S \times S},
\end{equation}

with entries given by the corresponding Pearson correlation values.

To obtain a sparse and interpretable representation graph,
an undirected weighted graph \(G_r=(S,E_r)\) was constructed by
thresholding the similarity matrix. Nodes in \(S\) represent cell lines.
Edges were retained when the pairwise similarity was greater than or equal to the empirical
90th percentile of off-diagonal pairwise cell-line similarities to define a
sparse weighted graph. Edge weights were set to the corresponding
Pearson similarity values.

For cell-line graphs, node degree was defined as the number of retained
incident edges. Betweenness centrality was computed on unweighted shortest
paths within each connected component. These definitions were used for
representation graphs, support-threshold consensus networks, the resolved
neighbour graph, and the pan-cancer transcriptomic similarity network where
applicable.

Community structure for representation-specific patient-referenced graphs was
evaluated with \texttt{tidygraph} Louvain clustering
(\texttt{group\_louvain}, using the \texttt{similarity} edge weight)
\citep{Blondel2008} and Leiden clustering (\texttt{group\_leiden}, using the
\texttt{similarity} edge weight and resolution \(=1.0\)) \citep{Traag2019}.
Leiden assignments were used for the exported community tables, and
Louvain--Leiden contingency tables were exported for comparison. Nodes with
zero degree were classified as outliers.

\subsubsection{Support-threshold consensus network across representations}

For each representation graph
\(G_r=(S,E_r)\), the retained edge set \(E_r\) was used to quantify
edge support across representations. For each unordered cell-line pair
\(e=(s_i,s_j)\), support was defined as
\(a(e)=|\{r\in\mathcal{R}:e\in E_r\}|\), with support fraction
\(f(e)=a(e)/|\mathcal{R}|\). A support-threshold consensus network
\(G_{\mathrm{support}}(m)\) was formed by retaining edges with
\(a(e)\geq m\):

\begin{equation}
G_{\mathrm{support}}(m)
=
(S,E_{\mathrm{support}}),
\quad
E_{\mathrm{support}}
=
\{e\in \bigcup_{r\in\mathcal{R}}E_r : a(e)\geq m\}.
\end{equation}

This network represented recurrent cell-line similarity relationships across
representation graphs. The descriptive support-network figures used the union
network \(m=1\), with edge support counts retained as annotations; support from
one representation and support from at least two representations were displayed
separately in the corresponding figures.


\subsubsection{Graph-based resolved neighbour graph}


Graph-based resolution used representation graphs to select final neighbour
relationships. Neighbours were identified using an intersection strategy that
integrated both cohort-wide and cell-line-specific optimal feature--distance
representations
(Algorithm~\ref{alg:graph_based_consensus_resolution_main};
Supplementary Algorithm~\ref{alg:supp_graph_based_consensus_resolution}). The resulting neighbourhoods
were stored as a resolved neighbour table.

A global-best feature--distance representation \(r^{\ast}\) was selected by
ranking representations on aggregate consensus metrics. Representations were
ordered primarily by the mean per-cell-line fraction of tumour associations
exceeding the consensus threshold, followed by the median consensus probability,
mean consensus probability, and representation identifier for deterministic
tie-breaking. This ranking defined the cohort-wide representation used to extract
the global-neighbour set.

\begin{equation}
    r^{\ast} =
    \underset{r \in \mathcal{R}}{\arg\max}
    \left(
    \mathrm{frac\_ge\_thr}(r),
    \tilde{p}(r)
    \right).
    \label{eq:global-best}
\end{equation}

where \(\mathrm{frac\_ge\_thr}(r)\) denotes the mean across cell lines of the
per-cell-line fraction of tumour associations exceeding the consensus threshold
for representation \(r\), and \(\tilde{p}(r)\) denotes the median across cell
lines of the per-cell-line median consensus probability. Remaining ties were
resolved by mean consensus probability and then representation identifier.

For a given representation \(r\), let \(N_r(i)\) denote the set
of neighbours of cell line $i$ in the cell line similarity graph
constructed under $r$, treating all edges as undirected.

For each cell line, the feature--distance representation with the highest
fraction of tumour associations exceeding the consensus threshold was
identified. If multiple representations were co-optimal for a cell line, all
tied representations were retained in \(T_i\). The local neighbour set was then
defined as the union of neighbours across these tied representations before
intersecting with the global-neighbour set.

\begin{equation}
    N_{\mathrm{global}}(i) =
    N_{r^{\ast}}(i),
    \label{eq:global-neighbours}
\end{equation}

\begin{equation}
    N_{\mathrm{local}}(i) =
    \bigcup_{r \in T_i}
    N_r(i),
    \label{eq:local-neighbours}
\end{equation}

The final set of stable neighbours was defined by intersection:

\begin{equation}
    N_{\mathrm{final}}(i) =
    N_{\mathrm{global}}(i)
    \cap
    N_{\mathrm{local}}(i).
    \label{eq:stable-neighbours}
\end{equation}

This intersection strategy restricted neighbour assignments to pairs supported
by both the cohort-wide global-best representation and the cell-line-specific
local-best representation set.

The resolved neighbour table was subsequently represented as an undirected
resolved neighbour graph \(G_{\mathrm{final}}\) for later annotation and
visualisation. The graph was decomposed into connected components, and
nodes with degree zero were retained as isolates. Within each non-isolate
component, the most-connected node was identified by maximum degree.
For node $v$,

\begin{equation}
    \mathrm{BC}(v) =
    \sum_{s \neq v \neq t}
    \frac{\sigma_{st}(v)}{\sigma_{st}},
    \label{eq:betweenness}
\end{equation}

where $\sigma_{st}$ is the total number of shortest paths between
nodes $s$ and $t$, and $\sigma_{st}(v)$ is the number of those paths
passing through $v$.
Nodes with the highest within-component betweenness centrality were annotated
as bridge-like component anchors.
Ties in centrality-based annotation were resolved deterministically by node identifier in the implementation. Isolates were not assigned bridge-like component anchors.

\begin{algorithm}[!ht]
\caption{Graph-based resolution of patient-referenced cell-line neighbourhoods}
\label{alg:graph_based_consensus_resolution_main}
\small
\begin{algorithmic}[1]
\Require Representation graphs \(G_r=(S,E_r)\), representation-level metrics, and per-cell-line consensus metrics.
\Ensure Resolved graph \(G_{\mathrm{final}}\), connected components, isolates, most-connected nodes, and bridge-like anchors.
\State Select global-best representation \(r^\ast\) by ranking aggregate consensus metrics: primarily the mean per-cell-line fraction of tumour associations exceeding the consensus threshold, followed by median consensus probability, mean consensus probability, and representation identifier.
\ForAll{cell lines \(i\)}
    \State Select local-best representation set \(T_i\) using the implemented per cell line primary criterion, fraction above the consensus threshold.
    \State Retain all exact co-optimal local ties in \(T_i\).
    \State Set \(N_{\mathrm{global}}(i)\gets N_{r^\ast}(i)\).
    \State Set \(N_{\mathrm{local}}(i)\gets\bigcup_{r\in T_i}N_r(i)\).
    \State Set \(N_{\mathrm{final}}(i)\gets N_{\mathrm{global}}(i)\cap N_{\mathrm{local}}(i)\).
\EndFor
\State Construct \(G_{\mathrm{final}}\) as an undirected graph from all stable-neighbour pairs.
\State Decompose \(G_{\mathrm{final}}\) into connected components and mark nodes with degree zero as isolates.
\ForAll{non-isolate connected components}
    \State Compute degree and betweenness centrality on unweighted shortest paths.
    \State Select the most-connected node by maximum degree, resolving ties deterministically by node identifier.
    \State Select the bridge-like anchor by maximum unweighted betweenness centrality, resolving ties deterministically by node identifier.
\EndFor
\State \Return resolved edge table, node annotation table, components, isolates, most-connected nodes, and bridge-like anchors.
\end{algorithmic}
\end{algorithm}

\subsubsection{Marker-gene prioritisation with DESeq2}

Marker-gene prioritisation was performed using \texttt{DESeq2} on raw integer
count matrices for cell lines selected from the resolved graph structure.
\texttt{DESeq2} modelled counts with a negative binomial generalised linear
model and estimated size factors, dispersions, log$_2$ fold changes, Wald
statistics, and Benjamini--Hochberg adjusted $p$-values. Size factors were
estimated from the lineage-specific count matrix before contrast-specific
models were fitted.

Contrasts were defined separately within each disease lineage after graph-based
resolution. For each contrast, a binary grouping variable was constructed and
modelled with design \(\sim\mathrm{group}\). The focal group was the tested isolate or
bridge-like anchor. The reference group was denoted ``REST'' for isolate
contrasts and ``OUTSIDE\_COMP'' for bridge-like anchor contrasts in the analysis
tables. Two complementary contrast strategies were applied:

\begin{enumerate}
	    \item \textbf{Isolate vs REST.}
	    Each graph isolate was contrasted against all remaining cell lines in the
	    same disease lineage. For isolate contrasts, REST therefore consisted of all
	    non-focal cell lines in that lineage.

	    \item \textbf{Bridge-like anchor vs REST.}
	    For each non-isolate resolved component, the bridge-like anchor was
	    contrasted against cell lines outside its assigned component. Cell lines in
	    the same resolved component as the anchor were excluded from REST so that the
	    comparator did not include the anchor's immediate graph neighbours.
\end{enumerate}

Replicate libraries were handled according to the profile identifiers present
in the lineage-specific raw count and graph tables. When a focal cell line had
more than one retained profile, all profiles assigned to that focal identifier
contributed to the focal group. Many focal contrasts nevertheless had only one
focal profile and a larger pooled comparator group.
These contrasts were therefore used to prioritise marker genes for graph-derived
cell-line states, not as strong population-level differential-expression
inference for low-replicate focal groups.

Genes were retained as marker genes if they satisfied the following criteria:

\begin{itemize}
    \item Benjamini–Hochberg adjusted $p$-value
    $\leq 0.01$ for isolates and
    $\leq 0.05$ for bridge-like anchors,
    \item mean normalised expression
    ($\mathrm{baseMean}$) $\geq 10$,
    \item a minimum absolute $\log_{2}$ fold change of $\geq 1.5$ for isolate contrasts
    and $\geq 1.0$ for bridge-like anchor contrasts,
    \item at least ten normalised counts in the tested focal sample.
\end{itemize}

Together, these criteria defined the expression-filtered marker-prioritisation
analysis. The adjusted \(p\)-value thresholds controlled multiple testing within
each DESeq2 contrast, whereas the fold-change, expression, and top-\(N\)
criteria were used to obtain interpretable marker lists for recurrence filtering
and enrichment analysis.

Marker lists were ranked by decreasing absolute \(\log_2\) fold change,
then by adjusted \(p\)-value and Ensembl gene identifier for deterministic
tie-breaking. Ranked lists were truncated to the top 50 genes per isolate
contrast and the top 200 genes per bridge-like anchor contrast.

Recurrence was subsequently quantified across contrasts (Supplementary
Algorithm~\ref{alg:supp_marker_recurrence_pan_cancer_features}). Within each
disease lineage, genes that passed the marker filters in at least two contrasts
were retained as strict recurrence-filtered genes. The pan-cancer feature set
was then extended by adding up to 50 additional isolate-versus-REST marker genes
per disease lineage that were not already in the recurrent core. Isolate-rescue
genes were ranked deterministically by isolate-contrast support, best adjusted
$p$-value, maximum absolute $\log_{2}$ fold change, marker-list rank, and
Ensembl gene identifier. The pan-cancer feature list was generated after
stripping Ensembl version suffixes and removing duplicate Ensembl identifiers
while preserving the first occurrence under the feature-selection order.

The resulting deduplicated pan-cancer feature set was used for the pan-cancer
cell-line similarity network and for the two correlation-based
similarity-ranking analyses described below. Because this feature set was
derived from graph-resolved marker-prioritisation contrasts within the study,
the ranking analyses were treated as internal transcriptomic similarity
analyses rather than external validation experiments.

\subsubsection{Pan-cancer transcriptomic similarity network of cell lines}

For cross-lineage visualisation, an undirected similarity graph using only
cell-line profiles was constructed from the pan-cancer feature set (Supplementary
Algorithm~\ref{alg:supp_pan_cancer_cell_line_network}).
The pan-cancer cell-line similarity network was constructed over the 167
analysed cell-line RNA-seq profiles defined in the dataset table. RBL replicate
libraries were retained as profile-level nodes in this network.
Each cell line profile was connected to its \(k = 20\) nearest neighbours by
Spearman correlation distance using a union \(k\)-nearest-neighbour graph; an
undirected edge was retained when either endpoint selected the other profile as
a nearest neighbour. Exact nearest-neighbour ties were resolved by profile
identifier, and edge weights were set to Spearman correlation values. Community
detection used \texttt{igraph::cluster\_louvain()} with edge weights and no
explicit Louvain resolution parameter for the reported partition;
\texttt{igraph::cluster\_leiden()} was also run with edge weights, resolution
\(=1.0\), and random seed 1 as an algorithmic comparison output.

Community structure was quantified by community size, within-community edge
count, mean within-community edge weight, and degree measures. Mean
within-community edge weight was defined as the arithmetic mean of Spearman
correlations across edges whose endpoints belonged to the same community.
The community purity statistic was defined as the fraction of profiles in a community assigned
to its majority lineage. A lineage-discordant profile was defined as a profile
whose annotated lineage differed from the majority lineage of its assigned
community.

Network modularity was computed as weighted Newman--Girvan modularity using
edge weights. Lineage assortativity was computed as unweighted nominal
assortativity of annotated lineage labels over the graph topology.

\subsubsection{Functional enrichment of marker sets}

Functional enrichment analysis was performed on marker gene sets using
g:Profiler through \texttt{gprofiler2} for organism \texttt{hsapiens} and
annotation release \texttt{e114\_eg62\_p19}. Input gene sets included
per-contrast marker lists, upregulated and downregulated subsets, disease-level
recurrent marker sets, the strict pan-cancer recurrent marker union, the final
pan-cancer feature set, and the isolate-rescued subset. The strict recurrent
marker union was defined as the union of genes passing strict recurrence
filtering across disease lineages. The isolate-rescued subset was defined as
genes present in the pan-cancer feature set but absent from the strict recurrent
marker union.

For per-contrast enrichment tests, the background universe was defined as the set
of genes included in the corresponding DESeq2 contrast after normalised-expression
filtering. For recurrent, pan-cancer, and isolate-rescue gene sets, the
background universe was the corresponding union of genes eligible for marker
selection in the contributing contrasts. Enrichment was evaluated for GO:BP,
GO:MF, GO:CC, KEGG, Reactome, WikiPathways, the g:Profiler transcription factor
binding source (TF; reported by g:Profiler as \texttt{Transfac}), Human Protein
Atlas (HPA), and CORUM sources using the configured g:Profiler archive. Multiple
testing was controlled using the g:SCS correction, and terms with
g:SCS-corrected \(p \leq 0.05\) were retained as enriched terms.

For each enriched term, the exported enrichment table retained the corrected
g:Profiler \(p\)-value, the input-set--term overlap size
(\texttt{intersection\_size}), input-set size, term size, effective domain
size, precision, recall, and enrichment ratio. Precision was defined as
\(\texttt{intersection\_size}\) divided by input-set size, recall as
\(\texttt{intersection\_size}/\texttt{term\_size}\), and enrichment
ratio as the observed input-set precision divided by the expected term
frequency in the effective domain.

For visualisation, enriched terms were reduced to a compact display set
by retaining predefined priority terms and top-scoring terms within each
displayed gene set category. Heatmap intensity was defined as
\(-\log_{10}\) g:SCS-corrected \(p\)-value capped at 30.

\subsection{Correlation-based transcriptomic similarity ranking and ranking-based classification}

Two complementary correlation-based ranking analyses were performed in the
curated pan-cancer feature space, defined by the final deduplicated pan-cancer
feature set: ranking cell-line profiles from each patient tumour and ranking
patient tumours after profile-to-group aggregation. Both analyses were
performed independently of CPI; CPI was not used as an input feature, weight, or
ranking criterion for the final correlation rankings.

\subsubsection{Ranking cell lines from each patient tumour}

To rank cell-line profiles by transcriptomic similarity, a correlation-based
ranking from patient tumours to cell lines was implemented (Supplementary
Algorithm~\ref{alg:supp_tumour_to_cell_line_mapping}).
The task was formulated as a transcriptomic similarity-ranking problem, with
ranking-based lineage classification derived from the top-ranked cell-line
match. Expression similarity between patient tumour sample $t$
and cell line $c$ was computed for all tumour--cell-line pairs in the configured ranking universe.
Patient-to-cell-line ranking used the 58 solid-tumour cell-line profiles
defined in the dataset table (29 breast cancer, 18 neuroblastoma, and 11
retinoblastoma). RBL\_15 and RBL\_20 each contributed two replicate libraries,
which were retained as separate profile-level models in this direction.
Before similarity scoring, expression values were gene-wise \(z\)-standardised separately within the tumour and cell-line cohorts.
For patient-to-cell-line ranking, Spearman correlation was computed between each retained patient tumour and each cell-line profile in the curated feature space.
For patient tumour sample $t$, all cell lines were ranked in descending order of $s(t,c)$.
Let
\[
c_{(1)} \in \arg\max_{c} \, s(t,c)
\]
denote a highest ranked cell line, with ties resolved by lexicographically ascending cell line identifier.
When reporting rank distributions across samples (for example, empirical cumulative distributions of ranks), tied ranks were averaged.

For each cell-line profile \(c\), an empirical cumulative distribution function
of ranks was computed as the fraction of patient tumours for which \(c\) ranked
at or below rank \(r\). For each patient tumour cohort, non-focal lineage models
were defined as all ranked cell-line models whose annotated cancer lineage
differed from the focal tumour lineage. For breast cancer tumours,
neuroblastoma and retinoblastoma models formed the non-focal comparator
distribution. For neuroblastoma tumours, breast cancer and retinoblastoma models
formed the comparator. For retinoblastoma tumours, breast cancer and
neuroblastoma models formed the comparator. At each rank, the median and
interquartile range across non-focal model ECDFs were used as an internal
negative-reference distribution rather than as a single manually selected
negative-control cell line. For visualisation, the four focal-lineage profiles
with the lowest median rank within each tumour lineage were displayed. Ties
were resolved by higher median similarity score and then by profile identifier.
For displayed top-10 fractions, two-sided 95\% confidence intervals for
binomial proportions were computed using the Wilson score interval.

The predicted tumour lineage was defined as
\[
\hat{y}_t = \mathrm{lin}\!\left(c_{(1)}\right),
\]
i.e., the lineage of the nearest cell line in transcriptomic space.

Ranking-based classification performance was evaluated by assigning each tumour to the lineage of its top-ranked cell-line match and comparing this predicted lineage with the known tumour lineage. Three similarity-ranking classification metrics
were computed across $T$ tumours.

Top-1 accuracy quantified the fraction of tumours whose predicted lineage matched the true tumour lineage,
\[
\mathrm{Acc@1} \;=\; \frac{1}{T}\sum_{t=1}^{T}
\mathbb{I}\!\left[\hat{y}_t=\mathrm{lin}(t)\right].
\]

Top-$k$ accuracy was evaluated at \(k=10\) and measured whether at least one correct-lineage cell line appeared among the top $k$ ranked candidates,
\[
\mathrm{Acc@k} \;=\; \frac{1}{T}\sum_{t=1}^{T}
\mathbb{I}\!\left[\exists\, r\le k:\mathrm{lin}\!\left(c_{(r)}\right)=\mathrm{lin}(t)\right].
\]

Mean Reciprocal Rank (MRR) quantified how early the first correct-lineage hit occurred,
\[
\mathrm{MRR} \;=\; \frac{1}{T}\sum_{t=1}^{T}
\begin{cases}
\frac{1}{r^\ast_t}, & r^\ast_t=\min\{r : \mathrm{lin}(c_{(r)})=\mathrm{lin}(t)\} \\
0, & \text{otherwise}.
\end{cases}
\]
MRR was computed on the full ranked cell line list; MRR@10 was computed separately as a bounded secondary metric when reported.

To characterise ranking-based classification confidence, two complementary heuristics were computed using \(k=10\).
Score separation (confidence delta) quantified the margin between the best match and the $k$th-ranked match,
\[
\Delta_t \;=\; s\!\left(t,c_{(1)}\right) - s\!\left(t,c_{(k')}\right),
\qquad
k'=\min(k,|\mathcal{L}|),
\]
where $|\mathcal{L}|$ denotes the number of available cell lines. Top-$k$ lineage purity measured the fraction
of top-$k$ candidates belonging to the tumour lineage,
\[
\pi_t \;=\; \frac{1}{k'}\sum_{r=1}^{k'}
\mathbb{I}\!\left[\mathrm{lin}\!\left(c_{(r)}\right)=\mathrm{lin}(t)\right].
\]

\subsubsection{Cell-line-centred tumour-class similarity and Precision@k}

To quantify tumour-class similarity from each analysed cell-line group, a
correlation-based ranking was performed in the curated pan-cancer feature space
(Supplementary Algorithm~\ref{alg:supp_cellline_tumour_class_similarity}).
Only patient tumour samples passing the ESTIMATE tumour-purity threshold of
\(\geq 0.7\) were retained. Spearman correlation was first computed between
each retained cell-line profile and each retained tumour sample. The final
ranking tables contained all retained solid-tumour samples that passed the
tumour-purity threshold. For ranking patient tumours from each biological
cell-line group, replicate libraries belonging to the same cell-line identity
were collapsed by taking the median tumour similarity score for each tumour
before ranking. Because breast cancer and neuroblastoma cell lines were each
represented by single libraries, only retinoblastoma replicate libraries were
collapsed; RBL\_15 and RBL\_20 each contributed two libraries that were grouped
into single biological cell-line identifiers, yielding 9 retinoblastoma groups.
The resulting 29 breast cancer, 18 neuroblastoma, and 9 retinoblastoma
identifiers defined 56 biological cell-line groups used as the primary units
for cell-line-to-tumour ranking.

For each biological cell-line group \(s\), retained patient tumour samples were
ranked in decreasing order of the group-level tumour similarity score
\(\rho(s,t)\), producing the ordered list
\[
\left( t_{(1)}, t_{(2)}, \ldots, t_{(T)} \right).
\]
Rank~1 denotes the tumour sample with the highest similarity to cell-line group
\(s\).

The cell-line-to-tumour ranking output was defined with one row per biological
cell-line group and tumour sample. The table reports tumour rank,
Spearman similarity score, rank percentile, tumour lineage, cell-line
lineage, and same-lineage status. Rank percentile was defined as tumour rank
divided by the number of retained tumour samples.

The lineage of the top-ranked tumour was compared to the labelled cell-line
lineage. Accuracy of the top-ranked tumour lineage was defined as the fraction
of biological cell-line groups for which the top-ranked tumour shared the
labelled cell-line lineage,
\[
\mathrm{Acc}_{\mathrm{top}} \;=\; \frac{1}{S}\sum_{s=1}^{S} \mathbb{I}\!\left[\mathrm{lin}\!\left(t_{(1)}^{s}\right)=\mathrm{lin}(s)\right],
\]
where \(S\) is the number of biological cell-line groups. Balanced accuracy was computed
as the unweighted mean of per-lineage accuracies. Two-sided 95\% confidence
intervals for proportions were computed using the Wilson score interval.
The top-ranked lineage confusion matrix tabulated labelled cell-line lineage
against the lineage of the top-ranked tumour, with cells reported as counts and
row percentages.

For a biological cell-line group \(s\), Precision@k was defined as
\[
\mathrm{Precision@}k(s) =
\frac{1}{k}\sum_{r=1}^{k}
\mathbb{I}\!\left[\mathrm{lin}\!\left(t_{(r)}^{s}\right)=\mathrm{lin}(s)\right].
\]
Mean Precision@k at each displayed value of \(k\) was calculated across
biological cell-line groups within each lineage. Lineage prevalence null expectations
were defined as the fraction of retained tumours belonging to the corresponding
cell-line lineage. The resulting curves were lineage-stratified Precision@k
curves. The Precision@k ribbons were 95\% percentile bootstrap
intervals over biological cell-line groups, using 2{,}000 bootstrap resamples.
Lineage-specific curves were bootstrapped within lineage. No random seed was
explicitly set for this bootstrap procedure in the original figure-generation
script. Consequently, exact bootstrap ribbon limits may vary slightly on rerun,
whereas deterministic rankings and point estimates are unchanged.

Same-lineage tumour ranking was quantified using the median rank of
same-lineage tumours and the fraction of same-lineage tumours matched within
the top 10\% of the ranked tumour list.

Per-group confidence was characterised by the Spearman correlation score of the top-ranked tumour, by the rank~1 minus rank~2 score margin, and by the rank~1 minus rank~10 score margin.

To quantify how high-ranking tumour samples were distributed across patient
consensus-network components, the top 50 ranked tumour samples for
each biological cell-line group were mapped to consensus network components
derived from the patient cohort. For each component $C_k$, the component size $n_k = |C_k|$ was defined
as the number of patient tumour samples assigned to $C_k$, and the component
purity was defined as
\[
\pi_k \;=\; \max_{\ell \in \mathcal{L}}\; \frac{|\{ t \in C_k : \mathrm{lin}(t) = \ell \}|}{n_k},
\]
i.e.\ the fraction of tumours within $C_k$ belonging to the most prevalent
lineage. Each component was annotated by its dominant tumour lineage. Component
composition for each biological cell-line group was reported as the fraction of
the top 50 tumours in each component.

For retinoblastoma cell lines represented by multiple libraries, rank measures
for profiles from replicate libraries were inspected separately as a
descriptive sensitivity check on the replicate-collapse step. This auxiliary
inspection was reported alongside, but was not used to recompute, the primary
group-level ranking metrics.

\subsection{Computational reproducibility and stochastic procedures}

The Snakemake workflow, configuration files, and custom scripts record the
analysis parameters used for the reported runs. Standalone k-means clustering
and PCA-based k-means wrappers set the random seed from the pipeline
configuration, with default seed 42, before evaluating \(k\) and using 20 random
initialisations. \texttt{ConsensusClusterPlus} was also called with a seed from
the same configuration, with default seed 42, for its 1{,}000 resampling
iterations.

UMAP was used for qualitative visualisation rather than quantitative inference.
The UMAP scripts used for tumour-neighbourhood quality-control plots and
DEG-set visualisations set a seed in code, either seed 42 or the configured seed.
A legacy pan-cancer UMAP helper did not contain an explicit seed. Because UMAP
coordinates were not used as inputs to clustering, graph construction,
marker-gene prioritisation, enrichment, or similarity ranking, stochastic
variation in UMAP affects visual coordinates rather than reported ranking or
graph statistics.

For graph procedures, deterministic nearest-neighbour sorting and tie-breaking
were used where specified. Pan-cancer Leiden community detection was run with
random seed 1 for the comparison output, whereas the representation-specific
\texttt{tidygraph} Leiden annotations did not pass an explicit seed in the
method call. Figure layout scripts used fixed archived coordinates or fixed
layout seeds where available; graph layouts were not used to calculate graph
statistics.

The cell-line-centred Precision@k ribbons used 2{,}000 percentile-bootstrap
resamples of biological cell-line groups. The original figure-generation script
did not set an explicit random seed for this bootstrap. Exact bootstrap ribbon
limits may therefore vary slightly on rerun, whereas the similarity scores,
rankings, Precision@k point estimates, top-ranked lineage calls, and Wilson
intervals are deterministic given the same input tables.

\subsection{Computational workflow}

\begin{figure}[H]
    \centering
    \includegraphics[height=0.58\textheight,keepaspectratio]{Thesis_writeup/Figures/workflow_diagram.pdf}
    \caption{Pan-cancer cell-line to patient tumour transcriptomic similarity workflow. RNA-seq profiles from the analysed cell-line panel and patient tumour cohorts entered the pipeline as independent inputs and were processed through sequencing quality control and alignment, tumour purity filtering, source harmonisation, feature selection, consensus clustering, tumour neighbourhood analysis, patient-referenced similarity modelling, graph-based consensus resolution, and similarity-ranking analyses. Colours indicate the major phases of the workflow.}
    \label{fig:workflow_diagram}
\end{figure}

High performance computing clusters at the Leibniz Institute DSMZ and the University of Potsdam were used to facilitate the analysis. Software dependencies were managed with Conda environments to support computational reproducibility, and the RNA-seq processing workflow was implemented using Snakemake.

\paragraph{Code availability.}
The Snakemake workflow, custom R and Python scripts, analysis configuration
files, and derived data needed to reproduce the reported analyses are available
at \url{https://github.com/EltonUgbogu/classification_of_dsmz_celllines.git}.
